% Estilo de la presentación del TP5

\usepackage[T1]{fontenc}
\usepackage[utf8]{inputenc}
\usepackage[spanish,es-noshorthands,es-nodecimaldot]{babel}
\spanishplainpercent

\usepackage{lmodern}
\usefonttheme{professionalfonts}
\usepackage{mathptmx}
\usepackage{amsmath,amssymb}
\DeclareSymbolFont{symbols}{OMS}{cmsy}{m}{n}
\SetSymbolFont{symbols}{bold}{OMS}{cmsy}{b}{n}
\DeclareSymbolFontAlphabet{\mathcal}{symbols}

\usepackage{graphicx}
\usepackage{tikz}
\usepackage{media9}
\usepackage{algpseudocode}
\algrenewcommand\algorithmicrequire{\textbf{Entrada:}}
\algrenewcommand\algorithmicensure{\textbf{Salida:}}
\newcommand{\acc}{a}   % a(r,v) = (-k r - \gamma v)/m

\usetikzlibrary{arrows.meta,positioning}
\graphicspath{{images/}}


\usetheme{Warsaw}
\setbeamertemplate{navigation symbols}{}
\definecolor{acento}{rgb}{0.2,0.2,0.7}
\definecolor{dlat}{HTML}{DB2B2B}   % mismos colores que D_LAT_COLOR / D_ACT_COLOR en scripts/lib/plot_style.py
\definecolor{dact}{HTML}{6E00B7}
\newcommand{\dlat}{\textcolor{dlat}{d_{\mathrm{lat}}}}   % los dos parámetros, siempre en su color (modo matemático)
\newcommand{\dact}{\textcolor{dact}{d_{\mathrm{act}}}}

% Barra de navegación: una columna por sección, un punto por diapositiva.
\useoutertheme[subsection=false]{miniframes}

% Sin esto las diapositivas salen sin numerar.
\setbeamertemplate{page number in head/foot}[framenumber]

% Separador con el título de la sección, sin numerar.
\AtBeginSection[]{%
  \begin{frame}[plain,noframenumbering]
    \vfill\centering\usebeamercolor[fg]{structure}%
    {\Huge\bfseries\insertsectionhead}\vfill
  \end{frame}%
}


\newcommand{\vv}[1]{\mathbf{#1}}
\newcommand{\ten}[1]{\times 10^{#1}}
\newcommand{\uni}[1]{\,\mathrm{#1}}
\newcommand{\va}{v_{a}}

\newcommand{\siimagen}[3]{%
  \IfFileExists{images/#1}{#2}{#3}%
}

\newcommand{\figura}[2][\linewidth]{%
  \siimagen{#2}%
    {\makebox[\linewidth][c]{\includegraphics[width=#1,height=0.74\textheight,keepaspectratio]{#2}}}%
    {\makebox[\linewidth][c]{\framebox[#1][c]{\rule{0pt}{5cm}\textit{falta la figura}}}}%
}


\newenvironment{cols}{\begin{columns}[c]}{\end{columns}}

\newenvironment{resultado}[1][Resultado]{%
  \begin{beamercolorbox}[wd=\linewidth,sep=6pt,rounded=true]{block body alerted}%
  \footnotesize\textbf{#1}\par\vspace{2pt}%
  \begin{itemize}\setlength{\itemsep}{1pt}%
}{%
  \end{itemize}\end{beamercolorbox}%
}

\newenvironment{parametros}[1][Parámetros]{%
  \begin{beamercolorbox}[wd=\linewidth,sep=6pt,rounded=true]{block body}%
  \footnotesize\textbf{#1}\par\vspace{2pt}%
  \begin{itemize}\setlength{\itemsep}{1pt}%
}{%
  \end{itemize}\end{beamercolorbox}%
}

\newenvironment{aclaraciones}[1][Aclaraciones]{%
  \begin{beamercolorbox}[wd=\linewidth,sep=6pt,rounded=true]{block body}%
  \footnotesize\textbf{#1}\par\vspace{2pt}%
  \begin{itemize}\setlength{\itemsep}{1pt}%
}{%
  \end{itemize}\end{beamercolorbox}%
}

\providecommand{\animbuild}{0}
\newif\ifanim
\ifnum\animbuild=1 \animtrue \else \animfalse \fi
\ifanim\usepackage{animate}\fi

\newcommand{\animacion}[6]{%
  \ifanim
    \animategraphics[autoplay,loop,width=\linewidth]{#4}{#1}{#2}{#3}%
  \else
    \siimagen{#5}%
      {\makebox[\linewidth][c]{\includegraphics[width=\linewidth,height=0.88\textheight,keepaspectratio]{#5}}}%
      {\makebox[\linewidth][c]{\framebox[\linewidth][c]{\rule{0pt}{4.2cm}\textit{fotograma}}}}%

    \if\relax\detokenize{#6}\relax\else
      \par\vspace{2pt}{\scriptsize\color{acento}\url{#6}}%
    \fi%
  \fi
}


% Logo ITBA al pie de la portada
\addtobeamertemplate{title page}{}{%
  \begin{tikzpicture}[remember picture,overlay]
    \node[
      anchor=south,
      yshift=0.35cm
    ] at (current page.south) {%
      \includegraphics[height=0.65cm]{images/Itba_bsas_logo.png}
    };
  \end{tikzpicture}%
}